\section{Related Work}

\textbf{Observer effects and Heisenbugs.} Heisenbugs and probe effects describe programs whose behavior changes under observation or debugging. Musuvathi et al. discuss debugging statements affecting concurrent interleavings in systematic testing \citep{musuvathi2008chess}. Sallinger, Weissenbacher, and Zuleger formalize Heisenbugs as hyperproperties and analyze causes including probe effects and nondeterminism \citep{sallinger2026heisenbugs}. OSDS is narrower: it models sequential access-state mechanisms where observation-shaped operations mutate latent state and later reads expose the difference.

\textbf{Semantic preservation and equivalence.} Compiler verification and translation validation study semantic preservation between source and transformed programs \citep{pnueli1998translation,leroy2009compcert}. OSDS uses the same preservation motivation but focuses on generated source-level transformations over library objects whose read and observation operations carry latent state. The oracle compares access-order behavior rather than proving full program equivalence.

\textbf{Metamorphic testing.} Metamorphic testing addresses oracle problems by checking relations across related executions \citep{chen1998metamorphic,chen2018metamorphic}. The OSDS-aware oracle is a specialized metamorphic relation: the follow-up execution changes the placement of an observation-shaped access while keeping witness setup matched. Generic metamorphic testing supplies the testing pattern; OSDS supplies the mechanism and classification boundary.

\textbf{Current generated-transformation benchmarks.} Differential fuzzing has recently been applied to LLM-generated code refactorings, showing that generated refactors can be functionally nonequivalent and that existing tests can miss a substantial fraction of those cases \citep{dristi2026differential}. SWE-Refactor moves refactoring evaluation to repository-level Java changes mined from real projects, emphasizing realistic context and compound refactorings \citep{xu2026swerefactor}. VERINA benchmarks verifiable code generation in Lean across code, specification, and proof generation tasks \citep{ye2025verina}. Viverra studies text-to-code with generated assertions and bounded verification, surfacing verified properties alongside generated C programs \citep{wu2026viverra}. Metamorphic transformations have also been used to diagnose memorization in LLM-based program repair; we use that line only as related evidence that semantics-preserving transformations can support stronger code-evaluation protocols \citep{dekoning2026metamorphic}.

\textbf{OSDS mechanism attribution.} Generic differential and metamorphic testing can establish that two executions differ under a chosen relation. OSDS contributes a sequential semantic model for objects whose reads and observations interact through latent access state. That structure determines which access reorderings form meaningful witnesses and supports mechanism-specific controls. The causal interventions test whether a detected divergence disappears when the access-induced mechanism is neutralized while the generated candidate remains unchanged, giving mechanism attribution alongside behavioral detection.

\textbf{Generated code and repair verification.} HumanEval and Codex made execution-based evaluation central for code generation \citep{chen2021codex}. SWE-bench extends evaluation to issue-resolution patches in real repositories \citep{jimenez2024swebench}. Automated program repair has long shown that patches can overfit tests \citep{legoues2012genprog,smith2015overfitting}. This paper studies a different preservation failure: a generated transformation can pass the ordinary benchmark tests while violating an access-state relation that those tests do not exercise.
